\begin{verse}
(1) The $y$-coordinates of two of the singularities coincide then
change order while the $x$-coordinates remain distinct.

(2) The $y$-coordinate along some string has a point of
inflection and the picture, before and after, looks locally like
one of the following
\end{verse}
\[
        \begin{picture}(0,0)%
\epsfig{file=ps/pic92.ps}%
\end{picture}%
\setlength{\unitlength}{0.00041700in}%
%
\begingroup\makeatletter\ifx\SetFigFont\undefined
% extract first six characters in \fmtname
\def\x#1#2#3#4#5#6#7\relax{\def\x{#1#2#3#4#5#6}}%
\expandafter\x\fmtname xxxxxx\relax \def\y{splain}%
\ifx\x\y   % LaTeX or SliTeX?
\gdef\SetFigFont#1#2#3{%
  \ifnum #1<17\tiny\else \ifnum #1<20\small\else
  \ifnum #1<24\normalsize\else \ifnum #1<29\large\else
  \ifnum #1<34\Large\else \ifnum #1<41\LARGE\else
     \huge\fi\fi\fi\fi\fi\fi
  \csname #3\endcsname}%
\else
\gdef\SetFigFont#1#2#3{\begingroup
  \count@#1\relax \ifnum 25<\count@\count@25\fi
  \def\x{\endgroup\@setsize\SetFigFont{#2pt}}%
  \expandafter\x
    \csname \romannumeral\the\count@ pt\expandafter\endcsname
    \csname @\romannumeral\the\count@ pt\endcsname
  \csname #3\endcsname}%
\fi
\fi\endgroup
\begin{picture}(10768,4496)(376,-4221)
\put(451,-436){\makebox(0,0)[lb]{\smash{\SetFigFont{12}{14.4}{rm}(A)}}}
\put(376,-3436){\makebox(0,0)[lb]{\smash{\SetFigFont{12}{14.4}{rm}(B)}}}
\put(6226,-511){\makebox(0,0)[lb]{\smash{\SetFigFont{12}{14.4}{rm}(C)}}}
\put(6226,-3511){\makebox(0,0)[lb]{\smash{\SetFigFont{12}{14.4}{rm}(D)}}}
\end{picture}

\]
In case 2), invariance of $Z_{\Theta}$ is guaranteed by (v) of 4.1.20
(in the order $C,D,B,A$) and in case 1), (i)$\to$(iv) of 4.1.20 is a
systematic enumeration of all 25 possibilities. For instance, the
case
\[
        \begin{picture}(0,0)%
\epsfig{file=ps/pic93.ps}%
\end{picture}%
\setlength{\unitlength}{0.00041700in}%
%
\begingroup\makeatletter\ifx\SetFigFont\undefined
% extract first six characters in \fmtname
\def\x#1#2#3#4#5#6#7\relax{\def\x{#1#2#3#4#5#6}}%
\expandafter\x\fmtname xxxxxx\relax \def\y{splain}%
\ifx\x\y   % LaTeX or SliTeX?
\gdef\SetFigFont#1#2#3{%
  \ifnum #1<17\tiny\else \ifnum #1<20\small\else
  \ifnum #1<24\normalsize\else \ifnum #1<29\large\else
  \ifnum #1<34\Large\else \ifnum #1<41\LARGE\else
     \huge\fi\fi\fi\fi\fi\fi
  \csname #3\endcsname}%
\else
\gdef\SetFigFont#1#2#3{\begingroup
  \count@#1\relax \ifnum 25<\count@\count@25\fi
  \def\x{\endgroup\@setsize\SetFigFont{#2pt}}%
  \expandafter\x
    \csname \romannumeral\the\count@ pt\expandafter\endcsname
    \csname @\romannumeral\the\count@ pt\endcsname
  \csname #3\endcsname}%
\fi
\fi\endgroup
\begin{picture}(10845,2146)(879,-3085)
\put(5401,-2011){\makebox(0,0)[lb]{\smash{\SetFigFont{12}{14.4}{rm}$\cdots$}}}
\put(11701,-2011){\makebox(0,0)[lb]{\smash{\SetFigFont{12}{14.4}{rm}$\cdots$}}}
\end{picture}

\]
\ni is covered by the first equation of (iii) of 4.1.20 with $j=i$.
Thus $Z_{\Theta}$ is invariant under isotopies preserving
standardness of the cusps.

Now we argue that a general isotopy $\phi_t$ may be replaced by a
$\tilde\phi_t$ for which $\tilde\phi_t$ preserves the standard form
of cusps for all $t$, without changing
$Z_{\Theta}$. To see this, construct a small disc around each cusp of
$\Theta$, sufficiently small so that the $y$-coordinates of all
points in a given disc are distinct from those in any other disc
and distinct from any maxima or minima of $y$ on the strings, and
such that the same is true for the images of these discs under
$\phi_1$. (Remember that $\phi_1(\Theta)$ is also a standard picture.)
Now in each disc $D$, construct a smaller disc $D_0$ inside $D$,
centered at the cusp, sufficiently small so that one can construct
a new isotopy $\tilde\phi_t$ having the properties

\begin{verse}
(i) $\tilde\phi_t$ restricted to each $D_0$ is just
translation in the plane\newline (and $\tilde\phi_t$ (a cusp)
$=\phi_t$ (that cusp)).

(ii) $\tilde\phi_t=\phi_t$ on the complement of the discs
$D$.
\end{verse}
Thus inside $\tilde\phi_t(D_0)$, the cusp remains standard and
$\tilde\phi_t$ is extended somehow to the annular region between
$\tilde\phi_t(D_0)$ and $\phi_t(D)$.  But the mapping class group
of diffeomorphisms of the annulus that are the identity on the
boundary is generated by a Dehn twist of 360${}^\circ$. So in a
neighborhood of each cusp point, $\phi_1$ and $\tilde\phi_1$ differ
only by some integer power of a single full twist. Figure~\ref{pic94}
illustrates  how $\phi_1(\Theta)$ and $\tilde\phi_1(\Theta)$ would
differ if the twist incurred were a single clockwise twist.
%%%%%%%%%%%%%%%%%%%%%%%%%%%%%%%%%%%%%%%%%%%%%%%%%%%%%%%%%%%%%%
\[
        \begin{picture}(0,0)%
\epsfig{file=ps/pic94.ps}%
\end{picture}%
\setlength{\unitlength}{0.00041700in}%
%
\begingroup\makeatletter\ifx\SetFigFont\undefined
% extract first six characters in \fmtname
\def\x#1#2#3#4#5#6#7\relax{\def\x{#1#2#3#4#5#6}}%
\expandafter\x\fmtname xxxxxx\relax \def\y{splain}%
\ifx\x\y   % LaTeX or SliTeX?
\gdef\SetFigFont#1#2#3{%
  \ifnum #1<17\tiny\else \ifnum #1<20\small\else
  \ifnum #1<24\normalsize\else \ifnum #1<29\large\else
  \ifnum #1<34\Large\else \ifnum #1<41\LARGE\else
     \huge\fi\fi\fi\fi\fi\fi
  \csname #3\endcsname}%
\else
\gdef\SetFigFont#1#2#3{\begingroup
  \count@#1\relax \ifnum 25<\count@\count@25\fi
  \def\x{\endgroup\@setsize\SetFigFont{#2pt}}%
  \expandafter\x
    \csname \romannumeral\the\count@ pt\expandafter\endcsname
    \csname @\romannumeral\the\count@ pt\endcsname
  \csname #3\endcsname}%
\fi
\fi\endgroup
\begin{picture}(9196,4811)(826,-4558)
\put(826,-4486){\makebox(0,0)[lb]{\smash{\SetFigFont{12}{14.4}{rm}$\phi_1(\Theta )$, near cusp}}}
\put(5776,-4486){\makebox(0,0)[lb]{\smash{\SetFigFont{12}{14.4}{rm}${\tilde\phi_1}(\Theta )$, near cusp.}}}
\end{picture}

\]
\begin{center}
	Figure 4.2.5
\end{center}
%%%%%%%%%%%%%%%%%%%%%%%%%%%%%%%%%%%%%%%%%%%%%%%%%
\bs We want to show, first in this case and then in the case of an
arbitrary integral power of a full twist, that $Z_{\Theta}$ is
unchanged. For this, consider Figure 4.2.6 which is supposed to be
part of a standard labelled picture in which the maxima and minima
of $y$, in the figure, and its cusps occur as an uninterrupted
sequence in
$\frak S(\Theta)$:
%%%%%%%%%%%%%%%%%%%%%%%%%%%%%%%%%%%%%%%%%%%%%%%%%%%%%%%%%
\[
        \begin{picture}(0,0)%
\epsfig{file=ps/pic95.ps}%
\end{picture}%
\setlength{\unitlength}{0.00041700in}%
%
\begingroup\makeatletter\ifx\SetFigFont\undefined
% extract first six characters in \fmtname
\def\x#1#2#3#4#5#6#7\relax{\def\x{#1#2#3#4#5#6}}%
\expandafter\x\fmtname xxxxxx\relax \def\y{splain}%
\ifx\x\y   % LaTeX or SliTeX?
\gdef\SetFigFont#1#2#3{%
  \ifnum #1<17\tiny\else \ifnum #1<20\small\else
  \ifnum #1<24\normalsize\else \ifnum #1<29\large\else
  \ifnum #1<34\Large\else \ifnum #1<41\LARGE\else
     \huge\fi\fi\fi\fi\fi\fi
  \csname #3\endcsname}%
\else
\gdef\SetFigFont#1#2#3{\begingroup
  \count@#1\relax \ifnum 25<\count@\count@25\fi
  \def\x{\endgroup\@setsize\SetFigFont{#2pt}}%
  \expandafter\x
    \csname \romannumeral\the\count@ pt\expandafter\endcsname
    \csname @\romannumeral\the\count@ pt\endcsname
  \csname #3\endcsname}%
\fi
\fi\endgroup
\begin{picture}(6097,3837)(2026,-4183)
\put(2026,-4111){\makebox(0,0)[lb]{\smash{\SetFigFont{12}{14.4}{rm}$y_-$}}}
\put(2026,-1036){\makebox(0,0)[lb]{\smash{\SetFigFont{12}{14.4}{rm}$y_+$}}}
\end{picture}

\]
\begin{center}
	Figure 4.2.6
\end{center}

\bs\ni From the definition of $Z(y)$ we see that
\begin{eqnarray*}
Z(y_+) & = & a_{{}_j}\mu_{{}_{j+1}} \a_{{}_{j+1,c}} \kappa_{{}_j}\eta_{{}_j} (Z(y_{{}_-})) \\ 
& = &\a_{{}_j,\rho(c)} (Z(y_{{}_-})) \quad \mbox{(by 4.1.13 and 4.1.14)}
\end{eqnarray*}
(Here the cusp is labelled by $c\in N'\cap M_{n-1}$ where
$n=3$ in Figure 4.2.6.) Similarly we see that if the cusp is
surrounded by a full 360${}^\circ$ twist,
\begin{eqnarray*}
Z(y_+) & = & \a_{j,\rho^n(b)} (Z(y_-)) \\
& = & \a_j, b (Z(y_-)) \qquad \mbox{(by Theorem  4.1.18)}
\end{eqnarray*}
Thus if a cusp  (or any part of a picture that is just a scaled
down standard labelled picture) is surrounded by a single clockwise
full twist, the effect on $Z_{\Theta}$ is as if the twist were not
there. If there were anticlockwise full twists around a cusp,
surround it further by the same number of clockwise twists. This
does not change
$Z_{\Theta}$, but the cancelling of the positive and negative
twists involves only isotopies that are the identity near the
cusps. Thus by our previous argument clockwise full twists around
cusps do not change
$Z_{\Theta}$ either. We conclude
\[
Z_{\phi_1(\Theta)} = Z_{\tilde\phi_1(\Theta)} =Z_{\Theta} \ .
\]

We have established that $Z_{\Theta}$ may be assigned to a standard
labelled picture by a product of elementary maps $\a, \kappa,
\eta,a,\mu$ in such a way that $Z_{\Theta}$ is unchanged by isotopies
of $\Theta$. We can now formally see how this makes $\coprod_n V_n$
a planar algebra according to $\S$1. The labelling set will be
$\coprod_n V_n$ itself.  The first step will be to associate a
labelled picture $\beta(T)$ with a labelled tangle $T$. To do this,
shrink all the internal boxes of $T$ to points and  isotope the
standard $k$-box to $[0,1]\x [0,1]$ with all the marked boundary
points going to points in $[0,1]\x 1$. Then distort the shrunk
boxes of $T$ to standard cusps, by isotopy, so that the string
attached to the first boundary point of the box becomes the first
string (from the left) attached to the cusp. The procedure near a
4-box of $T$ is illustrated in Figure~\ref{pic97} :
%%%%%%%%%%%%%%%%%%%%%%%%%%%%%%%%%%%%%%%%%%%%%%%%%%%%%%
\[
        \begin{picture}(0,0)%
\epsfig{file=ps/pic96.ps}%
\end{picture}%
\setlength{\unitlength}{0.00041700in}%
%
\begingroup\makeatletter\ifx\SetFigFont\undefined
% extract first six characters in \fmtname
\def\x#1#2#3#4#5#6#7\relax{\def\x{#1#2#3#4#5#6}}%
\expandafter\x\fmtname xxxxxx\relax \def\y{splain}%
\ifx\x\y   % LaTeX or SliTeX?
\gdef\SetFigFont#1#2#3{%
  \ifnum #1<17\tiny\else \ifnum #1<20\small\else
  \ifnum #1<24\normalsize\else \ifnum #1<29\large\else
  \ifnum #1<34\Large\else \ifnum #1<41\LARGE\else
     \huge\fi\fi\fi\fi\fi\fi
  \csname #3\endcsname}%
\else
\gdef\SetFigFont#1#2#3{\begingroup
  \count@#1\relax \ifnum 25<\count@\count@25\fi
  \def\x{\endgroup\@setsize\SetFigFont{#2pt}}%
  \expandafter\x
    \csname \romannumeral\the\count@ pt\expandafter\endcsname
    \csname @\romannumeral\the\count@ pt\endcsname
  \csname #3\endcsname}%
\fi
\fi\endgroup
\begin{picture}(11868,3880)(914,-3626)
\put(3428,-1522){\makebox(0,0)[lb]{\smash{\special{ps:gsave currentpoint currentpoint translate
-225.0 rotate neg exch neg exch translate}\SetFigFont{12}{14.4}{rm}1\special{ps:currentpoint grestore moveto}}}}
\put(3004,-1946){\makebox(0,0)[lb]{\smash{\special{ps:gsave currentpoint currentpoint translate
-225.0 rotate neg exch neg exch translate}\SetFigFont{12}{14.4}{rm}2\special{ps:currentpoint grestore moveto}}}}
\put(2580,-2370){\makebox(0,0)[lb]{\smash{\special{ps:gsave currentpoint currentpoint translate
-225.0 rotate neg exch neg exch translate}\SetFigFont{12}{14.4}{rm}3\special{ps:currentpoint grestore moveto}}}}
\put(2155,-2794){\makebox(0,0)[lb]{\smash{\special{ps:gsave currentpoint currentpoint translate
-225.0 rotate neg exch neg exch translate}\SetFigFont{12}{14.4}{rm}4\special{ps:currentpoint grestore moveto}}}}
\put(1943,-1734){\makebox(0,0)[lb]{\smash{\special{ps:gsave currentpoint currentpoint translate
-225.0 rotate neg exch neg exch translate}\SetFigFont{12}{14.4}{rm}6\special{ps:currentpoint grestore moveto}}}}
\put(1519,-2158){\makebox(0,0)[lb]{\smash{\special{ps:gsave currentpoint currentpoint translate
-225.0 rotate neg exch neg exch translate}\SetFigFont{12}{14.4}{rm}5\special{ps:currentpoint grestore moveto}}}}
\put(2368,-1310){\makebox(0,0)[lb]{\smash{\special{ps:gsave currentpoint currentpoint translate
-225.0 rotate neg exch neg exch translate}\SetFigFont{12}{14.4}{rm}7\special{ps:currentpoint grestore moveto}}}}
\put(2792,-885){\makebox(0,0)[lb]{\smash{\special{ps:gsave currentpoint currentpoint translate
-225.0 rotate neg exch neg exch translate}\SetFigFont{12}{14.4}{rm}8\special{ps:currentpoint grestore moveto}}}}
\put(2527,-1787){\makebox(0,0)[lb]{\smash{\special{ps:gsave currentpoint currentpoint translate
-225.0 rotate neg exch neg exch translate}\SetFigFont{12}{14.4}{rm}$R$\special{ps:currentpoint grestore moveto}}}}
\put(7501,-961){\makebox(0,0)[lb]{\smash{\SetFigFont{12}{14.4}{rm}1}}}
\put(7876,-1711){\makebox(0,0)[lb]{\smash{\SetFigFont{12}{14.4}{rm}2}}}
\put(7726,-2386){\makebox(0,0)[lb]{\smash{\SetFigFont{12}{14.4}{rm}3}}}
\put(7201,-2836){\makebox(0,0)[lb]{\smash{\SetFigFont{12}{14.4}{rm}4}}}
\put(5851,-2536){\makebox(0,0)[lb]{\smash{\SetFigFont{12}{14.4}{rm}5}}}
\put(5776,-1861){\makebox(0,0)[lb]{\smash{\SetFigFont{12}{14.4}{rm}6}}}
\put(6076,-1261){\makebox(0,0)[lb]{\smash{\SetFigFont{12}{14.4}{rm}7}}}
\put(6676,-1111){\makebox(0,0)[lb]{\smash{\SetFigFont{12}{14.4}{rm}8}}}
\put(10426,-2536){\makebox(0,0)[lb]{\smash{\SetFigFont{12}{14.4}{rm}1}}}
\put(10726,-2911){\makebox(0,0)[lb]{\smash{\SetFigFont{12}{14.4}{rm}$R$}}}
\end{picture}

\]
\begin{center}
	Figure 4.2.7
\end{center}
%%%%%%%%%%%%%%%%%%%%%%%%%%%%%%%%%%%%%%%%%%%%%%
\bs\ni The label associated to the cusp is just $\th^{-1}$ of the
label associated to the box but we may reasonably suppress
$\th^{-1}$.

Figure~4.2.8 illustrates a labelled tangle $T$ and a labelled
standard picture $\beta(T)$:
%%%%%%%%%%%%%%%%%%%%%%%%%%%%%%%%%%%%%%%%%%%%%%%%%%%%%%
\[
        \begin{picture}(0,0)%
\epsfig{file=ps/pic97.ps}%
\end{picture}%
\setlength{\unitlength}{0.00041700in}%
%
\begingroup\makeatletter\ifx\SetFigFont\undefined
% extract first six characters in \fmtname
\def\x#1#2#3#4#5#6#7\relax{\def\x{#1#2#3#4#5#6}}%
\expandafter\x\fmtname xxxxxx\relax \def\y{splain}%
\ifx\x\y   % LaTeX or SliTeX?
\gdef\SetFigFont#1#2#3{%
  \ifnum #1<17\tiny\else \ifnum #1<20\small\else
  \ifnum #1<24\normalsize\else \ifnum #1<29\large\else
  \ifnum #1<34\Large\else \ifnum #1<41\LARGE\else
     \huge\fi\fi\fi\fi\fi\fi
  \csname #3\endcsname}%
\else
\gdef\SetFigFont#1#2#3{\begingroup
  \count@#1\relax \ifnum 25<\count@\count@25\fi
  \def\x{\endgroup\@setsize\SetFigFont{#2pt}}%
  \expandafter\x
    \csname \romannumeral\the\count@ pt\expandafter\endcsname
    \csname @\romannumeral\the\count@ pt\endcsname
  \csname #3\endcsname}%
\fi
\fi\endgroup
\begin{picture}(8444,5194)(1779,-5833)
\put(3526,-1486){\makebox(0,0)[lb]{\smash{\special{ps:gsave currentpoint currentpoint translate
-180.0 rotate neg exch neg exch translate}\SetFigFont{12}{14.4}{rm}$R$\special{ps:currentpoint grestore moveto}}}}
\put(4201,-3811){\makebox(0,0)[lb]{\smash{\special{ps:gsave currentpoint currentpoint translate
-270.0 rotate neg exch neg exch translate}\SetFigFont{12}{14.4}{rm}$P$\special{ps:currentpoint grestore moveto}}}}
\put(3901,-2911){\makebox(0,0)[lb]{\smash{\special{ps:gsave currentpoint currentpoint translate
-180.0 rotate neg exch neg exch translate}\SetFigFont{12}{14.4}{rm}$Q$\special{ps:currentpoint grestore moveto}}}}
\put(7501,-2611){\makebox(0,0)[lb]{\smash{\SetFigFont{12}{14.4}{rm}$R$}}}
\put(9076,-2836){\makebox(0,0)[lb]{\smash{\SetFigFont{12}{14.4}{rm}$P$}}}
\put(8551,-4261){\makebox(0,0)[lb]{\smash{\SetFigFont{12}{14.4}{rm}$Q$}}}
\put(2626,-5761){\makebox(0,0)[lb]{\smash{\SetFigFont{12}{14.4}{rm}The tangle $T$}}}
\put(6451,-5761){\makebox(0,0)[lb]{\smash{\SetFigFont{12}{14.4}{rm}A standard picture $\beta (T)$}}}
\end{picture}

\]
\begin{center}
	Figure 4.2.8
\end{center}
%%%%%%%%%%%%%%%%%%%%%%%%%%%%%%%%%%%%%%%%%%%%%%

\bs\ni Note that $\beta(T)$ is not well defined, but two different
choices of $\beta(T)$ for a given $T$ will differ by an isotopy so
the map $\Phi$, \ $\Phi(T)=Z_{\beta(T)}$, gives a well defined
linear map from the universal planar algebra on $\coprod_k N'\cap
M_{k-1}$.

We now check that $\Phi$ makes $(P^{N\subset M})$ into a
connected spherical $C^*$-planar algebra. The first thing to check
is that
$\Phi$ is a homomorphism of filtered algebras. But if $T_1$ and
$T_2$ are labelled $k$-tangles, a choice of $\beta(T_1T_2)$ is
shown below $(k=3)$.
\[
        \begin{picture}(0,0)%
\epsfig{file=ps/pic98.ps}%
\end{picture}%
\setlength{\unitlength}{0.00041700in}%
%
\begingroup\makeatletter\ifx\SetFigFont\undefined
% extract first six characters in \fmtname
\def\x#1#2#3#4#5#6#7\relax{\def\x{#1#2#3#4#5#6}}%
\expandafter\x\fmtname xxxxxx\relax \def\y{splain}%
\ifx\x\y   % LaTeX or SliTeX?
\gdef\SetFigFont#1#2#3{%
  \ifnum #1<17\tiny\else \ifnum #1<20\small\else
  \ifnum #1<24\normalsize\else \ifnum #1<29\large\else
  \ifnum #1<34\Large\else \ifnum #1<41\LARGE\else
     \huge\fi\fi\fi\fi\fi\fi
  \csname #3\endcsname}%
\else
\gdef\SetFigFont#1#2#3{\begingroup
  \count@#1\relax \ifnum 25<\count@\count@25\fi
  \def\x{\endgroup\@setsize\SetFigFont{#2pt}}%
  \expandafter\x
    \csname \romannumeral\the\count@ pt\expandafter\endcsname
    \csname @\romannumeral\the\count@ pt\endcsname
  \csname #3\endcsname}%
\fi
\fi\endgroup
\begin{picture}(6772,3944)(451,-4583)
\put(1951,-3436){\makebox(0,0)[lb]{\smash{\SetFigFont{12}{14.4}{rm}$\beta (T_1)$}}}
\put(5401,-3661){\makebox(0,0)[lb]{\smash{\SetFigFont{12}{14.4}{rm}$\beta (T_2)$}}}
\put(451,-2611){\makebox(0,0)[lb]{\smash{\SetFigFont{12}{14.4}{rm}$y$}}}
\end{picture}

\]
\ni If $y$ is as marked, from the definition of $Z$,
$Z(y)=Z(\beta(T_1))\otimes Z(\beta(T_2))$. Moreover each pair of
maximal $\bnd$ \hskip8pt  contributes a factor $\chi_i$
$(=a_i\mu_{i+1})$ to $Z(y)$ as $y$ increases, so if $k$ is even,
part (1) of 4.1.24 gives $\Phi(T_1T_2)=\Phi(T_1)\Phi(T_2)$ and if
$k$ is odd (as in the figure) the last maximum has the black
region above so contribute a factor $\mu$ and part (ii) of 4.1.24
applies.

That $\Phi$ is compatible with the filtrations amounts to showing
that
\[
\begin{picture}(0,0)%
\epsfig{file=ps/pic110.ps}%
\end{picture}%
\setlength{\unitlength}{0.00041700in}%
%
\begingroup\makeatletter\ifx\SetFigFont\undefined
% extract first six characters in \fmtname
\def\x#1#2#3#4#5#6#7\relax{\def\x{#1#2#3#4#5#6}}%
\expandafter\x\fmtname xxxxxx\relax \def\y{splain}%
\ifx\x\y   % LaTeX or SliTeX?
\gdef\SetFigFont#1#2#3{%
  \ifnum #1<17\tiny\else \ifnum #1<20\small\else
  \ifnum #1<24\normalsize\else \ifnum #1<29\large\else
  \ifnum #1<34\Large\else \ifnum #1<41\LARGE\else
     \huge\fi\fi\fi\fi\fi\fi
  \csname #3\endcsname}%
\else
\gdef\SetFigFont#1#2#3{\begingroup
  \count@#1\relax \ifnum 25<\count@\count@25\fi
  \def\x{\endgroup\@setsize\SetFigFont{#2pt}}%
  \expandafter\x
    \csname \romannumeral\the\count@ pt\expandafter\endcsname
    \csname @\romannumeral\the\count@ pt\endcsname
  \csname #3\endcsname}%
\fi
\fi\endgroup
\begin{picture}(8217,1457)(1276,-2284)
\put(2026,-1241){\makebox(0,0)[lb]{\smash{\SetFigFont{12}{14.4}{rm}$R$}}}
\put(4052,-1242){\makebox(0,0)[lb]{\smash{\SetFigFont{12}{14.4}{rm}$R$}}}
\put(6361,-2191){\makebox(0,0)[lb]{\smash{\SetFigFont{12}{14.4}{rm}$R$}}}
\put(8701,-2221){\makebox(0,0)[lb]{\smash{\SetFigFont{12}{14.4}{rm}$R$}}}
\put(2701,-1246){\makebox(0,0)[lb]{\smash{\SetFigFont{12}{14.4}{rm}$) = \Phi ($}}}
\put(1276,-1231){\makebox(0,0)[lb]{\smash{\SetFigFont{12}{14.4}{rm}$\Phi ($}}}
\put(5053,-1261){\makebox(0,0)[lb]{\smash{\SetFigFont{12}{14.4}{rm}or $Z($}}}
\put(4603,-1261){\makebox(0,0)[lb]{\smash{\SetFigFont{12}{14.4}{rm}$)$}}}
\put(9493,-1336){\makebox(0,0)[lb]{\smash{\SetFigFont{12}{14.4}{rm}$)$}}}
\put(7228,-1291){\makebox(0,0)[lb]{\smash{\SetFigFont{12}{14.4}{rm}$)=Z($}}}
\end{picture}

\]
If $k$ is even this follows from 4.1.25 and if $k$ is odd it
follows from 4.1.26 (together with $\sum_{b\in B} bv_kb^*=\d
E_kE_{k-1}\dots E_2$  to take care of the factor $\kappa$
introduced by the minimum in the picture). So $\Phi$ is a
homomorphism of filtered algebras.

Annular invariance (and indeed the whole operadic picture) is easy.
If $T$ is an element of ${\cal P}(\coprod_k N'\cap M_{k-1})$ (linear
combination of tangles) with $\Phi(T)=0$, then if $T$ is surrounded
by an annular labelled tangle $A$, then we may choose
$\beta(\pi_A(T))$ to look like
\[
        \begin{picture}(0,0)%
\epsfig{file=ps/pic99.ps}%
\end{picture}%
\setlength{\unitlength}{0.00041700in}%
%
\begingroup\makeatletter\ifx\SetFigFont\undefined
% extract first six characters in \fmtname
\def\x#1#2#3#4#5#6#7\relax{\def\x{#1#2#3#4#5#6}}%
\expandafter\x\fmtname xxxxxx\relax \def\y{splain}%
\ifx\x\y   % LaTeX or SliTeX?
\gdef\SetFigFont#1#2#3{%
  \ifnum #1<17\tiny\else \ifnum #1<20\small\else
  \ifnum #1<24\normalsize\else \ifnum #1<29\large\else
  \ifnum #1<34\Large\else \ifnum #1<41\LARGE\else
     \huge\fi\fi\fi\fi\fi\fi
  \csname #3\endcsname}%
\else
\gdef\SetFigFont#1#2#3{\begingroup
  \count@#1\relax \ifnum 25<\count@\count@25\fi
  \def\x{\endgroup\@setsize\SetFigFont{#2pt}}%
  \expandafter\x
    \csname \romannumeral\the\count@ pt\expandafter\endcsname
    \csname @\romannumeral\the\count@ pt\endcsname
  \csname #3\endcsname}%
\fi
\fi\endgroup
\begin{picture}(5144,4844)(1179,-4883)
\put(3601,-2161){\makebox(0,0)[lb]{\smash{\SetFigFont{12}{14.4}{rm}$\cdots$}}}
\put(3301,-1186){\makebox(0,0)[lb]{\smash{\SetFigFont{12}{14.4}{rm}$\beta (A)$}}}
\put(3301,-2986){\makebox(0,0)[lb]{\smash{\SetFigFont{12}{14.4}{rm}$\beta (T)$}}}
\end{picture}

\]
\ni (Strictly speaking, one needs to consider such a picture for
each tangle in the linear combination forming $T$, and add.)
Clearly if $Z_{\beta(T)}=0$, so is $Z_{\beta(\pi_A(T))}$, since the
map $\a_{j,\th^{-1}(Z_{\beta(T)})}$ is applied in forming
$Z_{\beta(\pi_A(T))}$.

We now turn to planarity. By definition $V_0=\Bbb C$ so we only
need to show \ dim $V_{1,1}=1$.  A basis element of
${\cal P}_{1,1}(\coprod^\i_{k=1} N'\cap M_{k-1})$ is a 1-tangle $T$
with a vertical straight line and planar networks to the left and
right. We may choose $\beta(T)$ to be as depicted below
\[
        \begin{picture}(0,0)%
\epsfig{file=ps/pic100.ps}%
\end{picture}%
\setlength{\unitlength}{0.00041700in}%
%
\begingroup\makeatletter\ifx\SetFigFont\undefined
% extract first six characters in \fmtname
\def\x#1#2#3#4#5#6#7\relax{\def\x{#1#2#3#4#5#6}}%
\expandafter\x\fmtname xxxxxx\relax \def\y{splain}%
\ifx\x\y   % LaTeX or SliTeX?
\gdef\SetFigFont#1#2#3{%
  \ifnum #1<17\tiny\else \ifnum #1<20\small\else
  \ifnum #1<24\normalsize\else \ifnum #1<29\large\else
  \ifnum #1<34\Large\else \ifnum #1<41\LARGE\else
     \huge\fi\fi\fi\fi\fi\fi
  \csname #3\endcsname}%
\else
\gdef\SetFigFont#1#2#3{\begingroup
  \count@#1\relax \ifnum 25<\count@\count@25\fi
  \def\x{\endgroup\@setsize\SetFigFont{#2pt}}%
  \expandafter\x
    \csname \romannumeral\the\count@ pt\expandafter\endcsname
    \csname @\romannumeral\the\count@ pt\endcsname
  \csname #3\endcsname}%
\fi
\fi\endgroup
\begin{picture}(6044,3644)(1179,-3983)
\put(3226,-1561){\makebox(0,0)[lb]{\smash{\SetFigFont{12}{14.4}{rm}$T_1$}}}
\put(4801,-1561){\makebox(0,0)[lb]{\smash{\SetFigFont{12}{14.4}{rm}$T_2$}}}
\put(4051,-3511){\makebox(0,0)[lb]{\smash{\SetFigFont{12}{14.4}{rm}$T_3$}}}
\end{picture}

\]
\ni where there are 0-pictures inside
 the regions $T_1,T_2,\dots$. It is a simple consequence of our
formalism that a closed picture surrounded by a {\it white} region
simply contributes a scalar in a multiplicative way. This is
because one may first isotope the big picture so that all the
maxima and cusps in the 0-picture have $y$ coordinates in an
uninterrupted sequence in $\frak S$, and the last singularity must
be a maximum, shaded below, the first being a minimum, shaded
above. The final map will be an $a_j$ and will send the
contribution of the 0-picture to an element of $N'\cap N=\Bbb
C$. Thus we only need to see that 0-pictures inside a black
region contribute a scalar in a multiplicative way. But we may
isotope the big picture so that the singular $y$-values of the
0-picture occur in uninterrupted succession, and near the
0-picture the situation is as below:
\[
        \begin{picture}(0,0)%
\epsfig{file=ps/pic101.ps}%
\end{picture}%
\setlength{\unitlength}{0.00041700in}%
%
\begingroup\makeatletter\ifx\SetFigFont\undefined
% extract first six characters in \fmtname
\def\x#1#2#3#4#5#6#7\relax{\def\x{#1#2#3#4#5#6}}%
\expandafter\x\fmtname xxxxxx\relax \def\y{splain}%
\ifx\x\y   % LaTeX or SliTeX?
\gdef\SetFigFont#1#2#3{%
  \ifnum #1<17\tiny\else \ifnum #1<20\small\else
  \ifnum #1<24\normalsize\else \ifnum #1<29\large\else
  \ifnum #1<34\Large\else \ifnum #1<41\LARGE\else
     \huge\fi\fi\fi\fi\fi\fi
  \csname #3\endcsname}%
\else
\gdef\SetFigFont#1#2#3{\begingroup
  \count@#1\relax \ifnum 25<\count@\count@25\fi
  \def\x{\endgroup\@setsize\SetFigFont{#2pt}}%
  \expandafter\x
    \csname \romannumeral\the\count@ pt\expandafter\endcsname
    \csname @\romannumeral\the\count@ pt\endcsname
  \csname #3\endcsname}%
\fi
\fi\endgroup
\begin{picture}(5947,3644)(676,-3683)
\put(676,-1411){\makebox(0,0)[lb]{\smash{\SetFigFont{12}{14.4}{rm}$y_1$}}}
\put(676,-511){\makebox(0,0)[lb]{\smash{\SetFigFont{12}{14.4}{rm}$y_2$}}}
\put(3377,-2591){\makebox(0,0)[lb]{\smash{\SetFigFont{12}{14.4}{rm}picture}}}
\put(3377,-2087){\makebox(0,0)[lb]{\smash{\SetFigFont{12}{14.4}{rm}some}}}
\end{picture}

\]
\ni With $y_1$ as marked, $Z(y_1)$ will be $\sum_{b\in B}
b\otimes x\otimes b^*$ for some element $x$ in $N'\cap M$.
But then $Z(y_2)$ will be $\sum_{b\in B} bxb^*\in M'\cap M=\Bbb C$,
by 4.1.5.  So $P^{N\subset M}$ is a planar algebra.

The spherical property is easy: comparing
\bs
\[
        \begin{picture}(0,0)%
\epsfig{file=ps/pic102.ps}%
\end{picture}%
\setlength{\unitlength}{0.00041700in}%
%
\begingroup\makeatletter\ifx\SetFigFont\undefined
% extract first six characters in \fmtname
\def\x#1#2#3#4#5#6#7\relax{\def\x{#1#2#3#4#5#6}}%
\expandafter\x\fmtname xxxxxx\relax \def\y{splain}%
\ifx\x\y   % LaTeX or SliTeX?
\gdef\SetFigFont#1#2#3{%
  \ifnum #1<17\tiny\else \ifnum #1<20\small\else
  \ifnum #1<24\normalsize\else \ifnum #1<29\large\else
  \ifnum #1<34\Large\else \ifnum #1<41\LARGE\else
     \huge\fi\fi\fi\fi\fi\fi
  \csname #3\endcsname}%
\else
\gdef\SetFigFont#1#2#3{\begingroup
  \count@#1\relax \ifnum 25<\count@\count@25\fi
  \def\x{\endgroup\@setsize\SetFigFont{#2pt}}%
  \expandafter\x
    \csname \romannumeral\the\count@ pt\expandafter\endcsname
    \csname @\romannumeral\the\count@ pt\endcsname
  \csname #3\endcsname}%
\fi
\fi\endgroup
\begin{picture}(6547,2144)(2079,-3083)
\put(3376,-2086){\makebox(0,0)[lb]{\smash{\SetFigFont{12}{14.4}{rm}picture}}}
\put(6151,-2011){\makebox(0,0)[lb]{\smash{\SetFigFont{12}{14.4}{rm}and}}}
\put(7351,-2011){\makebox(0,0)[lb]{\smash{\SetFigFont{12}{14.4}{rm}picture}}}
\put(8626,-2086){\makebox(0,0)[lb]{\smash{\SetFigFont{12}{14.4}{rm},}}}
\end{picture}

\]
\bs\ni we see that the partition
functions are the same since one gives $E_{M'}$ applied to an
element of $N'\cap M$, and the other gives $E_N$ applied to the
same element with the correct powers of $\d$ contributed from the
minimum in the first picture and the maximum in the second. Either
way we get the trace by extremality.

For the $*$-structure, observe first that
$$
\th(x_1\otimes x_2\dots\otimes x_k)^* =
\th(x_k^*\otimes x_{k-1}^*\otimes\dots\otimes x_1^*)
$$
so
$\th(\a_{j,c}(x))^*=\th(\a_{k-j+2,c^*}(x^*))$, \
$\th(a_j(x))^*=\th(a_{k-j+1}(x^*))$, \
$\th(\mu_j(x))^*=\th(\mu_{n-j+2}(x^*))$, \
$\th(\eta_j(x))^*=\th(\eta_{k-j+2}(x^*))$ and
$\th(\kappa_j(x))^*=\th(\kappa_{k-j+1}(x^*))$.
Moreover if $T$ is a labelled tangle, $\beta(T^*)$ is $\beta(T)$
reflected in the line $x=\frac 12$ and with labels replaced by
their adjoints (via $\th$). We conclude that
$(Z_{\beta(T)})^*=Z_{\beta(T^*)}$ by applying the relations above
at each of the $y$-values in $\frak S(\beta(T))=\frak
S(\beta(T^*))$.

The $C^*$-property is just the positive definiteness of the
partition function. But if $T$ is a labelled $k$-tangle,
\[
\beta(\tangler{T}\ ) = \dtangle{\beta (T)} \ \ \ \ \mbox{so}
Z_{\Phi}(\tangler{T}\ )  =\d E_{M_{k-2}}(\Phi(\tangle{T}\ ))
\]
by 4.1.22.  Applying this $k$ times we get
\[
Z(\boxrloop{T} ) =
\d^k \ {\op{tr}}(\Phi(\tangle{T}\ ))
\]
so the positive definiteness of $Z$ follows from that of tr.

We now verify (i)$\to$(iv) in the statement of the theorem.

\smallskip
(i) Since $\Phi$ is a homomorphism of filtered algebras, it
suffices to prove the formula when $k=i$.
\begin{eqnarray*}
{\mbox{If}} \ k=2p+1, & \ Z_{\beta( | \ |\dots \quad )}= \th(Z_1( \butau )) \\
& = \d^{-P}\th\big( \sum_{b_1,b_2,\dots ,b_p} b_1\otimes b_2\otimes\dots \otimes b_p\otimes 1\otimes 1\otimes b^*_p \otimes\dots \otimes b^*_2\otimes b^*_1\big) \\
& = E_{2p+1} \qquad {\mbox{by Lemma 4.1.27.}}\vspace{2\jot}\\
\end{eqnarray*}
\begin{eqnarray*}
{\mbox {If  $k=2p$}},& \ Z_{\beta( | \ |\dots \quad )} =  \th(Z_1(\butad )) \\
&= \d^{-p-1}\th\big(\sum_{b_1\dots b_{p+1}} b_1\otimes b_2\otimes\dots \otimes b_p\otimes b^*_p b_{p+1}\otimes b^*_{p+1}\otimes b^*_{p-1} \otimes\dots \otimes b^*_1\big) \\
&= E_{2p} \qquad {\mbox{by  Corollary 4.1.28}}.
\end{eqnarray*}

(ii) The first formula follows from 4.1.5 and we showed the second
when we proved the positive defIniteness of the partition function.

(iii) This is just the filtered algebra property.

(iv) We also showed this in the  positive definiteness  proof.

All that remains is to prove the uniqueness of the planar algebra
structure. First observe that, as in Proposition 1.14, a labelled
tangle may be arranged by isotopy so that all of its boxes occur in
a vertical stack. After further isotopy and the introduction of
kinks or redundant loops, one may obtain the picture below for the
tangle (in ${\cal P}_k$)
\[
        \begin{picture}(0,0)%
\epsfig{file=ps/pic103.ps}%
\end{picture}%
\setlength{\unitlength}{0.00041700in}%
%
\begingroup\makeatletter\ifx\SetFigFont\undefined
% extract first six characters in \fmtname
\def\x#1#2#3#4#5#6#7\relax{\def\x{#1#2#3#4#5#6}}%
\expandafter\x\fmtname xxxxxx\relax \def\y{splain}%
\ifx\x\y   % LaTeX or SliTeX?
\gdef\SetFigFont#1#2#3{%
  \ifnum #1<17\tiny\else \ifnum #1<20\small\else
  \ifnum #1<24\normalsize\else \ifnum #1<29\large\else
  \ifnum #1<34\Large\else \ifnum #1<41\LARGE\else
     \huge\fi\fi\fi\fi\fi\fi
  \csname #3\endcsname}%
\else
\gdef\SetFigFont#1#2#3{\begingroup
  \count@#1\relax \ifnum 25<\count@\count@25\fi
  \def\x{\endgroup\@setsize\SetFigFont{#2pt}}%
  \expandafter\x
    \csname \romannumeral\the\count@ pt\expandafter\endcsname
    \csname @\romannumeral\the\count@ pt\endcsname
  \csname #3\endcsname}%
\fi
\fi\endgroup
\begin{picture}(3922,8131)(2101,-8054)
\put(4426,-6286){\makebox(0,0)[lb]{\smash{\SetFigFont{12}{14.4}{rm}$\vdots$}}}
\put(3751,-211){\makebox(0,0)[lb]{\smash{\SetFigFont{12}{14.4}{rm}$k$ boundary point}}}
\put(2101,-4036){\makebox(0,0)[lb]{\smash{\SetFigFont{12}{14.4}{rm}$T =$}}}
\put(4426,-2536){\makebox(0,0)[lb]{\smash{\SetFigFont{12}{14.4}{rm}$R_1$}}}
\put(4426,-4636){\makebox(0,0)[lb]{\smash{\SetFigFont{12}{14.4}{rm}$R_2$}}}
\put(4351,-1411){\makebox(0,0)[lb]{\smash{\SetFigFont{12}{14.4}{rm}$\sigma_1$}}}
\put(4351,-3511){\makebox(0,0)[lb]{\smash{\SetFigFont{12}{14.4}{rm}$\sigma_2$}}}
\put(4351,-7411){\makebox(0,0)[lb]{\smash{\SetFigFont{12}{14.4}{rm}$\sigma_n$}}}
\end{picture}

\]

\ni where the regions marked $\s_1,\dots ,\s_n$ contain only
strings and $\s_2,\dots ,\s_{n-1}$ have a fixed number $p$ of
boundary strings top and bottom, $\geq k$. Clearly $p-k$ is
even  so we conclude that, if $\Phi_1$ is some other planar
algebra structure satisfying (i) and (ii) then
$$
\Phi_1(T) =\d^{-\frac{p-k}{2}} \ \Phi_1
\big(\bigtan{T}\big) \  ,
$$
where we have introduced $\frac{p-k}{2}$ maxima and minima.
(To see this just apply the second formula of (ii) $p-k$ times.)
Thus we find that it suffices to prove that $\Phi'=\Phi$ on a
product of Temperley-Lieb tangles and tangles of the form
\[
        \begin{picture}(0,0)%
\epsfig{file=ps/pic104.ps}%
\end{picture}%
\setlength{\unitlength}{0.00041700in}%
%
\begingroup\makeatletter\ifx\SetFigFont\undefined
% extract first six characters in \fmtname
\def\x#1#2#3#4#5#6#7\relax{\def\x{#1#2#3#4#5#6}}%
\expandafter\x\fmtname xxxxxx\relax \def\y{splain}%
\ifx\x\y   % LaTeX or SliTeX?
\gdef\SetFigFont#1#2#3{%
  \ifnum #1<17\tiny\else \ifnum #1<20\small\else
  \ifnum #1<24\normalsize\else \ifnum #1<29\large\else
  \ifnum #1<34\Large\else \ifnum #1<41\LARGE\else
     \huge\fi\fi\fi\fi\fi\fi
  \csname #3\endcsname}%
\else
\gdef\SetFigFont#1#2#3{\begingroup
  \count@#1\relax \ifnum 25<\count@\count@25\fi
  \def\x{\endgroup\@setsize\SetFigFont{#2pt}}%
  \expandafter\x
    \csname \romannumeral\the\count@ pt\expandafter\endcsname
    \csname @\romannumeral\the\count@ pt\endcsname
  \csname #3\endcsname}%
\fi
\fi\endgroup
\begin{picture}(3944,944)(2379,-1283)
\put(4201,-961){\makebox(0,0)[lb]{\smash{\SetFigFont{12}{14.4}{rm}$x$}}}
\put(2701,-961){\makebox(0,0)[lb]{\smash{\SetFigFont{12}{14.4}{rm}$\cdots$}}}
\put(5401,-961){\makebox(0,0)[lb]{\smash{\SetFigFont{12}{14.4}{rm}$\cdots$}}}
\end{picture}

\]
\ni But the Temperley-Lieb algebra is known to be generated by
$\{\|\dots {}^{\cup}\!\!\!\!{}_{\cap_{i \ i+1}}\|\}$ whose images
are the $E_i$'s by (i). By condition (ii), we see that it suffices
to show that\newline $\Phi'(||\tangle{x})=
\Phi( || \tangle{x} )$.  
To this end we begin by showing
$$
\Phi'(\vln\vln\tanglel{a})= \Phi(\vln\vln\tanglel{a})
\quad {\text{for}} \ a\in N'\cap M_k \ .
$$
This follows from the picture below
\[
        \begin{picture}(0,0)%
\epsfig{file=ps/pic105.ps}%
\end{picture}%
\setlength{\unitlength}{0.00041700in}%
%
\begingroup\makeatletter\ifx\SetFigFont\undefined
% extract first six characters in \fmtname
\def\x#1#2#3#4#5#6#7\relax{\def\x{#1#2#3#4#5#6}}%
\expandafter\x\fmtname xxxxxx\relax \def\y{splain}%
\ifx\x\y   % LaTeX or SliTeX?
\gdef\SetFigFont#1#2#3{%
  \ifnum #1<17\tiny\else \ifnum #1<20\small\else
  \ifnum #1<24\normalsize\else \ifnum #1<29\large\else
  \ifnum #1<34\Large\else \ifnum #1<41\LARGE\else
     \huge\fi\fi\fi\fi\fi\fi
  \csname #3\endcsname}%
\else
\gdef\SetFigFont#1#2#3{\begingroup
  \count@#1\relax \ifnum 25<\count@\count@25\fi
  \def\x{\endgroup\@setsize\SetFigFont{#2pt}}%
  \expandafter\x
    \csname \romannumeral\the\count@ pt\expandafter\endcsname
    \csname @\romannumeral\the\count@ pt\endcsname
  \csname #3\endcsname}%
\fi
\fi\endgroup
\begin{picture}(2137,2055)(2529,-1726)
\put(3676,-811){\makebox(0,0)[lb]{\smash{\SetFigFont{12}{14.4}{rm}$a$}}}
\put(3451,-211){\makebox(0,0)[lb]{\smash{\SetFigFont{12}{14.4}{rm}$\cdots$}}}
\put(3451,-1411){\makebox(0,0)[lb]{\smash{\SetFigFont{12}{14.4}{rm}$\cdots$}}}
\end{picture}

\]
\ni for we know that
\[
\Phi'(|\tanglel{a})=
E_{M'}(a)=\Phi(|\tanglel{a}).
\]
\bs\ni
Now to show that \ $X=\Phi'( ||\tangle{x} )=
\Phi(||\tangle{x})=Y$  it suffices to
show that tr$(aX)={\op{tr}}(aY)$ for all
$a\in N'\cap M_k$. But up to powers of $\d$,
\[
	\begin{picture}(0,0)%
\epsfig{file=ps/pic111.ps}%
\end{picture}%
\setlength{\unitlength}{0.00041700in}%
%
\begingroup\makeatletter\ifx\SetFigFont\undefined
% extract first six characters in \fmtname
\def\x#1#2#3#4#5#6#7\relax{\def\x{#1#2#3#4#5#6}}%
\expandafter\x\fmtname xxxxxx\relax \def\y{splain}%
\ifx\x\y   % LaTeX or SliTeX?
\gdef\SetFigFont#1#2#3{%
  \ifnum #1<17\tiny\else \ifnum #1<20\small\else
  \ifnum #1<24\normalsize\else \ifnum #1<29\large\else
  \ifnum #1<34\Large\else \ifnum #1<41\LARGE\else
     \huge\fi\fi\fi\fi\fi\fi
  \csname #3\endcsname}%
\else
\gdef\SetFigFont#1#2#3{\begingroup
  \count@#1\relax \ifnum 25<\count@\count@25\fi
  \def\x{\endgroup\@setsize\SetFigFont{#2pt}}%
  \expandafter\x
    \csname \romannumeral\the\count@ pt\expandafter\endcsname
    \csname @\romannumeral\the\count@ pt\endcsname
  \csname #3\endcsname}%
\fi
\fi\endgroup
\begin{picture}(10545,2069)(751,-5108)
\put(3977,-3737){\makebox(0,0)[lb]{\smash{\SetFigFont{12}{14.4}{rm}$a$}}}
\put(4277,-4637){\makebox(0,0)[lb]{\smash{\SetFigFont{12}{14.4}{rm}$x$}}}
\put(7427,-4637){\makebox(0,0)[lb]{\smash{\SetFigFont{12}{14.4}{rm}$x$}}}
\put(7202,-3737){\makebox(0,0)[lb]{\smash{\SetFigFont{12}{14.4}{rm}$a$}}}
\put(10727,-4712){\makebox(0,0)[lb]{\smash{\SetFigFont{12}{14.4}{rm}$x$}}}
\put(10502,-3812){\makebox(0,0)[lb]{\smash{\SetFigFont{12}{14.4}{rm}$a$}}}
\put(2401,-3736){\makebox(0,0)[lb]{\smash{\SetFigFont{12}{14.4}{rm}tr($\Phi'($}}}
\put(4726,-3736){\makebox(0,0)[lb]{\smash{\SetFigFont{12}{14.4}{rm}$)) = $tr$(\Phi'($}}}
\put(7876,-3736){\makebox(0,0)[lb]{\smash{\SetFigFont{12}{14.4}{rm}$)) = $tr$(\Phi ($}}}
\put(11251,-3811){\makebox(0,0)[lb]{\smash{\SetFigFont{12}{14.4}{rm}$)) = $tr$(aY)$}}}
\put(751,-3736){\makebox(0,0)[lb]{\smash{\SetFigFont{12}{14.4}{rm}tr($aX)=$}}}
\end{picture}

\]
%%%%%%%%%
%\medskip
%\[
%{\op{tr}}(aX)  = {\op{tr}}\left(\Phi' \left( \qquad\qquad
%\right)\right) = {\op{tr}}\left(\Phi' \left( \qquad\qquad
%\right)\right)
% = {\op{tr}}\left(\Phi \left( \qquad\qquad
%\right)\right) = {\op{tr}}(aY) \ ,
%\]
%%%%%%%%%%%%%%%%%%%%%%%%%
\bs\ni
and we are done. \qed

\bs\ni{\bf Definition 4.2.8.} The {\it annular Temperley Lieb}
algebra $AT(n,\d)$, for $n$ even, will be the $*$-algebra with
presentation:
\[
        \begin{picture}(0,0)%
\epsfig{file=ps/pic106.ps}%
\end{picture}%
\setlength{\unitlength}{0.00041700in}%
%
\begingroup\makeatletter\ifx\SetFigFont\undefined
% extract first six characters in \fmtname
\def\x#1#2#3#4#5#6#7\relax{\def\x{#1#2#3#4#5#6}}%
\expandafter\x\fmtname xxxxxx\relax \def\y{splain}%
\ifx\x\y   % LaTeX or SliTeX?
\gdef\SetFigFont#1#2#3{%
  \ifnum #1<17\tiny\else \ifnum #1<20\small\else
  \ifnum #1<24\normalsize\else \ifnum #1<29\large\else
  \ifnum #1<34\Large\else \ifnum #1<41\LARGE\else
     \huge\fi\fi\fi\fi\fi\fi
  \csname #3\endcsname}%
\else
\gdef\SetFigFont#1#2#3{\begingroup
  \count@#1\relax \ifnum 25<\count@\count@25\fi
  \def\x{\endgroup\@setsize\SetFigFont{#2pt}}%
  \expandafter\x
    \csname \romannumeral\the\count@ pt\expandafter\endcsname
    \csname @\romannumeral\the\count@ pt\endcsname
  \csname #3\endcsname}%
\fi
\fi\endgroup
\begin{picture}(12473,4242)(1201,-4581)
\put(3676,-811){\makebox(0,0)[lb]{\smash{\SetFigFont{12}{14.4}{rm}$\cdots$}}}
\put(3676,-2011){\makebox(0,0)[lb]{\smash{\SetFigFont{12}{14.4}{rm}$\cdots$}}}
\put(5176,-1411){\makebox(0,0)[lb]{\smash{\SetFigFont{12}{14.4}{rm},}}}
\put(1351,-1411){\makebox(0,0)[lb]{\smash{\SetFigFont{12}{14.4}{rm}$F_1=$}}}
\put(5177,-3812){\makebox(0,0)[lb]{\smash{\SetFigFont{12}{14.4}{rm},}}}
\put(9527,-1412){\makebox(0,0)[lb]{\smash{\SetFigFont{12}{14.4}{rm},}}}
\put(9527,-3812){\makebox(0,0)[lb]{\smash{\SetFigFont{12}{14.4}{rm},}}}
\put(5702,-1412){\makebox(0,0)[lb]{\smash{\SetFigFont{12}{14.4}{rm}$F_2=$}}}
\put(9902,-1412){\makebox(0,0)[lb]{\smash{\SetFigFont{12}{14.4}{rm}$F_k=$}}}
\put(9902,-3812){\makebox(0,0)[lb]{\smash{\SetFigFont{12}{14.4}{rm}$F_{2k}=$}}}
\put(5476,-3811){\makebox(0,0)[lb]{\smash{\SetFigFont{12}{14.4}{rm}$F_{k+2}=$}}}
\put(1201,-3811){\makebox(0,0)[lb]{\smash{\SetFigFont{12}{14.4}{rm}$F_{k+1}=$}}}
\put(7651,-3211){\makebox(0,0)[lb]{\smash{\SetFigFont{12}{14.4}{rm}$\cdots$}}}
\put(7651,-4396){\makebox(0,0)[lb]{\smash{\SetFigFont{12}{14.4}{rm}$\cdots$}}}
\put(11866,-4336){\makebox(0,0)[lb]{\smash{\SetFigFont{12}{14.4}{rm}$\cdots$}}}
\put(11836,-3211){\makebox(0,0)[lb]{\smash{\SetFigFont{12}{14.4}{rm}$\cdots$}}}
\put(8042,-812){\makebox(0,0)[lb]{\smash{\SetFigFont{12}{14.4}{rm}$\cdots$}}}
\put(8042,-2012){\makebox(0,0)[lb]{\smash{\SetFigFont{12}{14.4}{rm}$\cdots$}}}
\put(12256,-811){\makebox(0,0)[lb]{\smash{\SetFigFont{12}{14.4}{rm}$\cdots$}}}
\put(12256,-2011){\makebox(0,0)[lb]{\smash{\SetFigFont{12}{14.4}{rm}$\cdots$}}}
\put(3256,-4411){\makebox(0,0)[lb]{\smash{\SetFigFont{12}{14.4}{rm}$\cdots$}}}
\put(3271,-3211){\makebox(0,0)[lb]{\smash{\SetFigFont{12}{14.4}{rm}$\cdots$}}}
\end{picture}

\]
\bs{\bf Remark.} Since $(F_1F_2\dots F_k)(F_1F_2\dots F_k)^*=\d F_1$
and $(F_1F_2\dots F_k)^*(F_1F_2\dots F_k)=\d F_k$, if $\cal H$
is a Hilbert space carrying a $*$-representation of $AT(n,\d)$,
dim$(F_i\cal H)$ is independent of $i$.

\bigskip
\ni{\bf Corollary 4.2.9} If $N\subset M$ is an extremal subfactor
of index $\d^{-2} >4$, each $N'\cap M_{k-1}$ is a Hilbert space
carrying a $*$-representation of $AT(2k,\d)$. And
${\op{dim}}(F_i(N'\cap M_{k-1}))={\op{dim}}(N'\cap M_{k-2})=
{\op{dim}}(M'\cap M_{k-1})$.

\bigskip
{\sc Proof.} Let $F_i$ be the elements of $\cal A(\phi)$ defined
as follows

\newpage

\ni The relations are easily checked, the Hilbert space structure
on $N'\cap M_{k-1}$ being given by the trace
$-\<a,b\>={\op{tr}}(b^*a)$. Since
$F_k(N'\cap M_{k-1})=N'\cap M_{k-2}$ and $F_{2k}
(N'\cap M_{k-1})=M'\cap M_{k-1}$ (by (ii) of 4.2.1).
We are through. \qed

\bigskip
\ni{\bf Lemma 4.2.10} If $\cal H$ carries an irreducible
 $*$-representation of $AT(n,\d)$ for $\d >2$, $n>4$, and
${\op{dim}}(F_i\cal H)=1$, then ${\op{dim}} \ \cal H=n$
(remember $n$ is even).

\bigskip
{\sc Proof.} Let $v_i$ be a unit vector in $F_i\cal H$ for each
$i$. Then $F_jv_i$ is a multiple of $v_j$ so the linear span of the
$v_i$'s is invariant, thus equal to $\cal H$ by irreducibility.
Here dim $\cal H\leq n$ (this does not require $\d >2$). Moreover
the commutation relations imply $|\<v_i,v_{i+1}\>|=\d^{-1}$ and
$\<v_i,v_j\>=0$ or 1 if $i\neq j\pm 1$. The case
$\<v_i,v_j\>=1$ forces $i=j$ so it only happens if $n=4$.
So by changing the $v_i$'s by phases we may assume that the matrix
$\d\<v_i,v_j\>$ is
\[
\D_n(\o) =\left(
\begin{array}{clcr}
\d & 1 & 0 & 0  \dots  \o \\
 1 & \d & 1 & 0  \dots  0 \\
 0 & 1 & \d & 1  \dots  0 \\
\vdots & {} & {} & {}    \vdots \\
\bar\o & 0 & \dots & {}   {1 \ \ \d}
\end{array}\right).
\]
It is easy to check that det$(\D_n(\o))=P_{2n}(\d)-P_{2n-2}(\d)-2
{\op{Re}}(\o)$ where $P_n(\d)$ are Tchebychev polynomials.
Thus det$(\D_n(\o))$ is smallest, for fixed $\d$, when $\o=1$.
But then \newline $\|\D_n(1)-\d \ {\op{id}}\|=2$ by Perron-Frobenius
so det $\D_n(1) >0$ for $\d >2$. \qed

\bs\ni
{\bf Corollary 4.2.11} Suppose the principal graph of the
subfactor $N\subset M$, $[M:N]>4$, has an initial segment equal to
the Coxeter-Dynkin diagram $D_{n+2}$ with $*$ as shown:
\[
        \begin{picture}(0,0)%
\epsfig{file=ps/pic107.ps}%
\end{picture}%
\setlength{\unitlength}{0.00041700in}%
%
\begingroup\makeatletter\ifx\SetFigFont\undefined
% extract first six characters in \fmtname
\def\x#1#2#3#4#5#6#7\relax{\def\x{#1#2#3#4#5#6}}%
\expandafter\x\fmtname xxxxxx\relax \def\y{splain}%
\ifx\x\y   % LaTeX or SliTeX?
\gdef\SetFigFont#1#2#3{%
  \ifnum #1<17\tiny\else \ifnum #1<20\small\else
  \ifnum #1<24\normalsize\else \ifnum #1<29\large\else
  \ifnum #1<34\Large\else \ifnum #1<41\LARGE\else
     \huge\fi\fi\fi\fi\fi\fi
  \csname #3\endcsname}%
\else
\gdef\SetFigFont#1#2#3{\begingroup
  \count@#1\relax \ifnum 25<\count@\count@25\fi
  \def\x{\endgroup\@setsize\SetFigFont{#2pt}}%
  \expandafter\x
    \csname \romannumeral\the\count@ pt\expandafter\endcsname
    \csname @\romannumeral\the\count@ pt\endcsname
  \csname #3\endcsname}%
\fi
\fi\endgroup
\begin{picture}(6172,1244)(826,-1283)
\put(5626,-436){\makebox(0,0)[lb]{\smash{\SetFigFont{12}{14.4}{rm}$n$}}}
\put(3826,-736){\makebox(0,0)[lb]{\smash{\SetFigFont{12}{14.4}{rm}$\cdots$}}}
\put(826,-736){\makebox(0,0)[lb]{\smash{\SetFigFont{12}{14.4}{rm}$*$}}}
\end{picture}

\]
\ni Then there are at least two edges of the principal graph
connecting the two points at distance $n+1$ from $*$ to points of
distance $n+2$.

\bs
{\sc Proof.} Since $\d >2$, the Temperley-Lieb algebra generated
by $\{1,e_1,e_2,\dots ,e_k\}$ in $N'\cap M_k$ has dimension
$\frac{1}{k+2}{2k+4 \choose k+1}$. The information on
the principal graph then gives the following Bratteli diagram for
the inclusions $N'\cap M_{n-1}\subset N'\cap M_n\subset
N'\cap M_{n+1}$:
\begin{eqnarray*}
N'&\cap & M_{n+1} \\
&\cup &\\
N'&\cap & M_n\\
&\cup&\\
N'&\cap & M_{n-1}
\end{eqnarray*}
Here we have shown only that part of the Bratteli diagram relevant
to the proof. Only the two 1's in the middle row can be connected
to anything in the top row other than vertices corresponding to the
ideal generated by $e_{n+1}$. We have to show that it is impossible
for just one of these 1's to be connected, with multiplicity one,
to a new principal graph vertex. By contradiction, suppose this
were the case. Then we would have
$$
{\op{dim}}(N'\cap M_{n+1})=\frac{1}{n+3}
{2(n+2)\choose n+2} -(n+1)^2+(n+2)^2
$$
since the only difference between the $N'\cap M_{n+1}$ level
of the Bratteli diagram and the Temperley-Lieb Bratteli diagram
(see [GHJ]) is that the ``$n+1$" in Temperley-Lieb has become
``$n+2$".  Thus
$$
{\op{dim}}(N'\cap M_{n+1})=\frac{1}{n+3}
{2(n+2)\choose n+2} +2n+3 \ .
$$
But consider $N'\cap M_{n+1}$ as a module over $AT(2n+4,\d)$.
The Temperley-Lieb subalgebra is invariant and so therefore is its
orthogonal complement $TL^{\perp}$ of dimension $2n+3$. But
consider the image of $F_{n+2}=\d E_{N'\cap M_n}$. It is
$\dsize{\frac{1}{n+2}{2(n+1)\choose n+1} +1}$ because of the
single extra vertex on the principal graph at distance $n+1$ from $*$.


